\section{2 октября 1914 г.}

\lp{20260604_193317492}

Вы спрашиваете, как мы нашли все, вернувшись из Крыма. Все было в порядке, и
квартира, и собаки.

\lp{20260604_193317492}

\noindent
У меня такое чувство, что это было страшно давно. Жили мы в Юдино недолго.
Только было Коля наладил свою работу и я свою работу, как
разразилась\textellipsis \rem{война}. Стали чуть ли не каждый день ездить в
город. Один раз, за неимением поездов для простых смертных, даже пошли пешком в
Кунцево. Это было в дни, когда брали Колю, и мы еще не знали, что его
освободят. А потом так стали волноваться, когда приходилось до вечера ждать
газет, что решили скорее переехать в город, и переехали 15 августа. Про
окружающее настроение что же писать? Конечно, совсем не то,что было в японскую
войну. Тогда все относились к ней, как к чему-то принудительному, затеянному
нами же, и даже не считали возможным желать России победы. Так, например,
относились такие люди как Киселев, Ильин. А сейчас все идут охотно, подъем
огромный, все убеждены, что нас вызвали, что Россия вовсе не хотела войны, и
что наше правительство не виновато и не оно принуждает идти на войну, а внешние
обстоятельства, и есть масса, масса

\lp{20260604_193311828}

\noindent
добровольцев. Даже старые, старые военные вроде отца Ольги Николаевны
Ковалевской просятся на войну. Подъем всеобщий. А я ночей не сплю и молю Бога
чтобы скорее кончилось. Вас я очень, очень люблю, мой родной, близкий! И когда
читала, что вы переживали и как вы слушали «Тристана и Изольду», то прямо
разрывалась. Но только помните, как вы сами соглашались с тем, что германцы
должны возродиться, что немцы сейчас в упадке. Это говорили часто и вы, и Коля,
когда возмущались на многие их теперешние проявления. И вот может быть потому
то еще, и в особенности так трагично и тяжело все это для вас, что в сущности
даже и в этом вашем положительном к ним отношении вы одиноки, т.к. они сами это
все (за очень вероятно небольшим исключением) не понимают что и как они должны
нести. И \rem*{зигерридства} своего не знают, забыли. Вообще, я все время делаю
усилие, чтобы в воображении своем слить любимые, знакомые черты, но отошедшие,
с теми, которые налицо. Невольно вспоминаю из настоящего те физиономии, которые
видала, и мужские, и женские и хочется забыть их и помнить другие, из прошлого,
которые в тясячу раз ближе, роднее и дороже; а слить их воедино никак не могу.
Вы и Коля из тех, прежних, старых и дорогих.

Ваша Анна.

Верочка тоже приходит в ваши комнатки. Играет иногда на пианино.
